\begin{table}[H]
\centering
\caption{Structural Robustness Using Assignment-Value Primitives}
\label{tab:struct_assignment_value_sensitivity}
\begin{threeparttable}
\footnotesize
\begin{tabular}[t]{lcccccc}
\toprule
Payoff primitive & Kenya & Liberia & Nigeria & Max ITT/SE & K high-$\tau$ & High-input \\
\midrule
Preferred rho & 0.533 & 0.055 & 0.121 & 0.281 & $+0.155$ & $+0.187$ \\
Net mismatch reduction & 0.533 & 0.320 & 0.000 & 0.282 & $+0.152$ & $+0.186$ \\
Share-weighted reduction & 0.533 & 0.238 & 0.000 & 0.282 & $+0.153$ & $+0.186$ \\
\bottomrule
\end{tabular}
\begin{tablenotes}[para,flushleft]
\footnotesize
\item \textit{Notes:} The preferred model uses the control-group predictive-content primitive $\rho$ in the assignment-payoff term. The robustness rows re-estimate the stage-4 production block after replacing that payoff input with a normalized assignment-value primitive from Table~\ref{tab:assignment_value_summary}. ``Net mismatch reduction'' uses the positive part of Grade minus Diagnostic mismatch, normalized so Kenya equals the preferred high-signal value. ``Share-weighted reduction'' multiplies the positive mean reduction by the share of students whose predicted mismatch falls before the same normalization. The mechanical sorting equation continues to use $\rho$ because it predicts classroom compression rather than payoff value. The table reports the Kenya high-delivery and fully high-input counterfactuals under each re-estimated production map.
\end{tablenotes}
\end{threeparttable}
\end{table}
